\documentclass[10pt]{article}

\usepackage[margin=1in]{geometry}
\usepackage[T1]{fontenc}
\usepackage{lmodern}
\usepackage{microtype}
\usepackage{amsmath,amssymb}
\usepackage{booktabs}
\usepackage{graphicx}
\usepackage{xcolor}
\usepackage{hyperref}
\usepackage{caption}

\definecolor{linkblue}{HTML}{1D4ED8}
\hypersetup{
  colorlinks=true,
  linkcolor=linkblue,
  citecolor=linkblue,
  urlcolor=linkblue,
  pdftitle={Simulating AR(1) and MA(1) Processes in R},
  pdfauthor={Kalle Saarelma}
}

\title{\textbf{Simulating AR(1) and MA(1) Processes in R}\\
\large A Reproducible Introduction to Linear Time-Series Dependence}
\author{Kalle Saarelma}
\date{July 2026}

\begin{document}

\maketitle

\begin{center}
\fcolorbox{black}{gray!8}{%
  \parbox{0.92\linewidth}{%
    \centering
    \textbf{AI-Generation Disclaimer}\\[0.35em]
    This article, including its prose, structure, analysis, and typesetting,
    was fully generated using artificial intelligence tools. It was created
    solely as an experiment to test the capabilities of AI-assisted research,
    coding, and document-production workflows. It should not be presented as
    independently authored or peer-reviewed academic research.
  }%
}
\end{center}
\vspace{0.6em}

\begin{abstract}
Autoregressive and moving-average processes provide the elementary building
blocks of linear time-series analysis. This article develops a compact,
reproducible implementation of first-order autoregressive, AR(1), and
first-order moving-average, MA(1), processes in base R. The implementation
separates simulation, visualization, and validation into reusable components
and records all parameter values and random seeds. Simulated trajectories and
sample autocorrelation functions illustrate the central theoretical
distinction between the models: the AR(1) autocorrelation decays geometrically,
whereas the MA(1) autocorrelation cuts off after lag one. The project
demonstrates how a small amount of well-structured code can connect formal
time-series theory with transparent computational experiments.
\end{abstract}

\noindent\textbf{Keywords:} autoregression, moving average, autocorrelation,
stationarity, reproducible research, R

\section{Introduction}

Time-series observations are ordered, and their temporal order is often
informative. Economic indicators, sensor measurements, financial returns, and
environmental records may all exhibit dependence between present and past
values. Linear stochastic processes offer an interpretable framework for
representing such dependence. In particular, autoregressive and moving-average
models form the basis of the ARMA and ARIMA families used throughout applied
forecasting and econometrics \cite{box2015,brockwell2016}.

This article studies the simplest nontrivial members of these families: the
AR(1) and MA(1) models. Although elementary, they display qualitatively
different dependence structures. An AR(1) process propagates earlier values
forward through a recursive equation. An MA(1) process, by contrast, depends
only on the current and immediately preceding shocks. Simulation makes this
contrast visible and provides a useful check on analytical results.

The accompanying R project has three aims. First, it implements both models
without external package dependencies. Second, it produces repeatable figures
and data by fixing random seeds. Third, it includes lightweight checks that
verify output dimensions and deterministic reproduction. The complete source
is available at
\url{https://github.com/kjsaarel/time-series-processes}.

\section{Model definitions and properties}

\subsection{The AR(1) process}

Let $\{\varepsilon_t\}$ be independent and identically distributed innovations
with mean zero and variance $\sigma_\varepsilon^2$. An AR(1) process satisfies
\begin{equation}
  y_t = \phi y_{t-1} + \varepsilon_t,
  \label{eq:ar1}
\end{equation}
where $\phi$ controls persistence. If $|\phi|<1$, the process has a
covariance-stationary solution with
\begin{equation}
  \operatorname{Var}(y_t)
  = \frac{\sigma_\varepsilon^2}{1-\phi^2}.
\end{equation}
Its theoretical autocorrelation function (ACF) is
\begin{equation}
  \rho(k)=\phi^k,\qquad k=0,1,2,\ldots.
  \label{eq:aracf}
\end{equation}
Positive values of $\phi$ produce persistence with a monotone geometric decay.
Negative values produce an alternating ACF. When $\phi=1$, Equation
\eqref{eq:ar1} becomes a random walk and is no longer covariance stationary.

\subsection{The MA(1) process}

An MA(1) process is defined by
\begin{equation}
  y_t = \varepsilon_t + \theta\varepsilon_{t-1},
  \label{eq:ma1}
\end{equation}
where $\theta$ determines the contribution of the previous innovation. The
variance and ACF are
\begin{align}
  \operatorname{Var}(y_t)
    &= (1+\theta^2)\sigma_\varepsilon^2,\\
  \rho(1)
    &= \frac{\theta}{1+\theta^2},\\
  \rho(k)
    &= 0,\qquad k\geq 2.
  \label{eq:maacf}
\end{align}
Thus, the MA(1) process is stationary for every finite $\theta$. The commonly
imposed condition $|\theta|<1$ concerns invertibility: it ensures a unique,
convergent autoregressive representation of the model \cite{hamilton1994}.

\section{Computational design}

The project separates the computational workflow into three layers. The file
\path{R/simulate_processes.R} contains functions for simulation and plotting.
The script \path{scripts/run_demo.R} fixes the experimental parameters,
generates both series, and writes figures and comma-separated data files.
Finally, \path{tests/test_simulations.R} checks vector lengths, initial
conditions, and reproducibility under fixed seeds.

Table \ref{tab:parameters} summarizes the baseline experiments. Innovations are
drawn from a standard normal distribution, so
$\sigma_\varepsilon^2=1$.

\begin{table}[htbp]
  \centering
  \caption{Simulation parameters.}
  \label{tab:parameters}
  \begin{tabular}{lcccc}
    \toprule
    Model & Observations & Parameter & Seed & Initial value \\
    \midrule
    AR(1) & 150 & $\phi=0.9$ & 12 & $y_1=0$ \\
    MA(1) & 100 & $\theta=0.5$ & 123 & Not applicable \\
    \bottomrule
  \end{tabular}
\end{table}

The AR(1) implementation uses an explicit recursion. This choice mirrors the
mathematical definition and makes the role of the initial condition
transparent. The MA(1) implementation draws $n+1$ innovations, including the
pre-sample shock $\varepsilon_0$, and then evaluates Equation
\eqref{eq:ma1} vectorially. Both simulation functions return the realized
series, innovations, and model parameter in a structured list.

Reproducibility requires more than setting a seed in an interactive session.
The parameters, seeds, output locations, and execution commands must also be
recorded. The complete experiment can be regenerated from the repository root
with
\begin{verbatim}
Rscript scripts/run_demo.R
Rscript tests/test_simulations.R
\end{verbatim}
No contributed R packages are required.

\section{Simulation results}

Figure \ref{fig:ar1} shows the AR(1) realization and its sample ACF. With
$\phi=0.9$, shocks have long-lived effects. The simulated trajectory therefore
exhibits extended excursions above and below zero rather than rapidly
oscillating around the mean. The sample ACF remains strongly positive at the
first several lags and declines gradually, in agreement with the geometric
pattern predicted by Equation \eqref{eq:aracf}. Sampling variation prevents the
empirical correlations from lying exactly on the theoretical sequence.

\begin{figure}[htbp]
  \centering
  \includegraphics[width=0.80\textwidth]
    {output/ar1_simulation.png}
  \caption{A simulated AR(1) process with $\phi=0.9$ (left) and its sample
  autocorrelation function (right).}
  \label{fig:ar1}
\end{figure}

Figure \ref{fig:ma1} displays the MA(1) experiment. Compared with the persistent
AR(1) trajectory, the MA(1) series reverses direction more rapidly because an
innovation affects only two adjacent observations. For $\theta=0.5$, the
theoretical first-lag correlation is
\[
  \rho(1)=\frac{0.5}{1+0.5^2}=0.4.
\]
The sample ACF shows a positive first-lag association and correlations near
zero thereafter. This abrupt cutoff is the diagnostic theoretical signature
of an MA(1) process.

\begin{figure}[htbp]
  \centering
  \includegraphics[width=0.80\textwidth]
    {output/ma1_simulation.png}
  \caption{A simulated MA(1) process with $\theta=0.5$ (left) and its sample
  autocorrelation function (right).}
  \label{fig:ma1}
\end{figure}

\section{Discussion}

The experiment illustrates how model structure determines serial dependence.
In an AR(1) model, a shock influences $y_t$ directly, $y_{t+1}$ through
$\phi$, $y_{t+2}$ through $\phi^2$, and so forth. The effect diminishes but does
not terminate after a fixed lag. In an MA(1) model, the same shock enters only
two consecutive observations. Consequently, the population ACF cuts off
exactly after lag one.

These observations also motivate a classical identification heuristic. An
autoregressive model tends to have a gradually decaying ACF and a partial
autocorrelation function that cuts off. A moving-average model tends to display
the reverse pattern \cite{shumway2017}. In finite samples, however, the
distinction is imperfect. Sample autocorrelations are noisy, and a highly
persistent AR process can resemble a nonstationary series. Model selection
should therefore combine graphical diagnostics with estimation, information
criteria, residual analysis, and substantive knowledge.

The present exercise is deliberately limited. It assumes Gaussian,
homoskedastic innovations; it does not estimate parameters from observed data;
and it considers only one realization per parameter setting. A natural
extension would compare empirical and theoretical ACFs across repeated
simulations. Further extensions could examine negative parameters,
near-unit-root behavior, forecast uncertainty, or mixed ARMA processes.

\section{Conclusion}

AR(1) and MA(1) processes encode two fundamentally different mechanisms of
time dependence. Autoregression transmits past values recursively and produces
geometrically decaying autocorrelation. Moving-average dependence transmits
recent innovations over a finite horizon and produces an ACF cutoff. The
reproducible R project developed here makes these properties visible while
maintaining a clear separation among simulation, output generation, and
validation. This combination of mathematical derivation and transparent code
provides a practical foundation for more advanced time-series modeling.

\small
\begin{thebibliography}{9}
\setlength{\itemsep}{0.25em}

\bibitem{box2015}
G. E. P. Box, G. M. Jenkins, G. C. Reinsel, and G. M. Ljung,
\textit{Time Series Analysis: Forecasting and Control}, 5th ed.
Hoboken, NJ: Wiley, 2015.

\bibitem{brockwell2016}
P. J. Brockwell and R. A. Davis,
\textit{Introduction to Time Series and Forecasting}, 3rd ed.
Cham: Springer, 2016.

\bibitem{hamilton1994}
J. D. Hamilton,
\textit{Time Series Analysis}.
Princeton, NJ: Princeton University Press, 1994.

\bibitem{shumway2017}
R. H. Shumway and D. S. Stoffer,
\textit{Time Series Analysis and Its Applications: With R Examples},
4th ed. Cham: Springer, 2017.

\end{thebibliography}

\end{document}
